\section{Forecasting task and evaluation protocol}\label{sec:protocol}

\paragraph{Task.}
Entity $e$ reports modal prices $P_{e,t}$ (INR per quintal) on irregular dates.
For an origin $t$ and $h\in\{3,7,14\}$ days, the target is the first report dated in $[t+h,t+h+3]$. If there is none, there is no label.
The \emph{next} horizon is the next report, if it falls within 7 days.
Models predict $y=\ln P_{e,\tau}-\ln P_{e,t}$, and errors are measured in rupees on $P_{e,t}\exp(\hat y)$.

\paragraph{Periods (main experiment).}
\begin{itemize}
  \item Training origins: 2016--2023.
  \item Validation: 2024, used for tuning, model choice and conformal calibration.
  \item Test origins: 2025-01-01 to 2026-04-20, scored once.
\end{itemize}
A row is kept only if its target date also lies in its period; a unit test enforces this purge.
A perturbation test multiplies all future prices, and all other entities' prices, by arbitrary factors and requires every feature at the origin to be unchanged.
Training uses a uniform sample of $2\times10^6$ rows per horizon (fixed seed).

\paragraph{Selection.}
The family with the lowest validation MAE is selected, baselines included.
The search space was amended twice, both times on validation evidence, before any test metric was read:
\begin{itemize}
  \item first, after squared-loss learners tied persistence;
  \item second, after finding that MLlib's absolute-loss boosting is mis-scaled for log returns.
\end{itemize}
An aborted first run had already written test predictions for one horizon. They were deleted unread, and we disclose this.

\paragraph{Metrics and tests.}
We report:
\begin{itemize}
  \item MAE, RMSE, WAPE, sMAPE and bias;
  \item MASE, scaled by each entity's in-sample $h$-step naive MAE on training rows \cite{hyndman2006another};
  \item MAE relative to persistence on identical rows;
  \item the share of series (at least 30 test rows) on which a model beats persistence.
\end{itemize}
Diebold--Mariano tests \cite{diebold1995comparing} use daily mean absolute-error differences, a Newey--West variance with lag $h+3$ (7 for the next report) and the HLN correction. Same-day errors are aggregated first because they are correlated across markets.

\paragraph{Frozen-model and retraining evaluations.}
After the official back-fill, we score the frozen models, quantiles and MASE scales on complete official data:
\begin{itemize}
  \item on the original test period;
  \item on 2026-04-22 to 2026-10-22, reports first downloaded after the models were saved.
\end{itemize}
Separately, we retrain with the same search space, selection rule, sample size and seed, on a protocol fixed before the run:
\begin{itemize}
  \item training: 2016--2024;
  \item validation: 2025;
  \item final test: 2026-01-01 to 2026-10-22, not used for any retraining choice. Its first months overlap the main test, and its later months had already been used to score the frozen models.
\end{itemize}


\paragraph{Locked retrospective robustness protocol.}
A separate resource-limited sensitivity experiment uses 2024 training origins, January--June 2025 for selection and early stopping, July--December 2025 for interval calibration, and January--October 2026 for retrospective evaluation. Targets are purged at each period boundary. It compares the original cleaned features with a second pipeline retaining reports flagged by the fivefold jump rule; other validity checks and panel membership are preserved. The second pipeline begins in October 2022, providing over a year of history before training origins. This is not an exact reproduction of the original feature construction, whose source state is unproven.
Each regime trains on 500,000 deterministic hash-selected rows, tests 31 and 63 leaves with at most 800 boosting rounds and one fixed seed, and compares all five simple baselines on identical rows. MASE scales use all training-period rows. Calibration observations have no role in hyperparameter choice or early stopping. Calendar-day paired moving-block bootstrap intervals use 14- and 28-day blocks and 500 replicates; separate entity-cluster intervals describe between-series sensitivity, not joint two-way inference. Eight primary candidate--persistence tests (two regimes, four horizons) are corrected together using Holm. For these tests, calendar-day total loss differences divided by a constant target pooled row-weighted MAE; the older daily-mean tests target a different weighting. These changes do not retrospectively correct the original protocol.
The protocol, imported module paths, source commit, dependencies, split bounds and SHA256 hashes of inputs, models and saved predictions are recorded. Generated tables are checked against saved predictions. All 2026 outcomes were already available during study design, so this experiment measures retrospective robustness and cannot constitute independent confirmation.


\paragraph{Release reproduction and dependence sensitivity.}
The clean-source release rebuilds both cleaning regimes from the same silver data and fixed panel, rather than reusing either historical feature table. Source, silver/panel inputs, shard definitions and the assembled executable JAR are hashed before execution, verified unchanged around shard execution, and bound to output checksums, commands, runtime versions and exit codes. This is observed execution attestation, not an independent build witness. It reproduces the smaller study, not the full-size historical MLlib experiments or the original panel-construction process.
A subsequent automated methodological review found that daily loss dependence can extend beyond the original bandwidths. Block choice matters under dependent resampling \cite{politis2004block}; we therefore add explicitly retrospective sensitivities with all 28-, 42- and 56-day blocks, circular and non-circular, and 2,000 replicates each. We report asymptotic normal HAC diagnostics at those three bandwidths without a fixed-horizon HLN correction. The older helper coupled that horizon correction to the covariance bandwidth; for irregular report gaps its finite-sample justification is unestablished. Those historical p values are retained only as heuristic diagnostics. Holm corrections are separate across each declared eight-comparison bandwidth family, with no favourable bandwidth selected.
Inference is conditional on the fitted models and observed panel, and relies on an adequate weak-dependence and stationarity approximation. Same-day aggregation retains joint observed cross-market variation, but does not eliminate serial dependence or seasonality. Longer blocks preserve more dependence while leaving fewer temporal spans; more replicates improve Monte Carlo resolution, not sampling validity. The interval estimates have not been shown to attain nominal coverage. Entity resampling is exploratory sensitivity under correlated series, not a confidence guarantee for independent new markets. We report direct entity-ratio quantiles, subgroup errors, daily losses/counts, autocorrelation and MASE exclusions. A human expert statistical review remains outstanding.
